\documentclass[tikz,border=1cm]{standalone}
\usepackage{lua-tikz3dtools}
\def\NU{5}
\def\NV{3}
\pgfmathtruncatemacro{\streamcount}{\NU*\NV}
\begin{document}
\foreach \i in {1,...,100} {
\pgfmathsetmacro{\i}{\i/100*0.99*2*pi}
\begin{tikzpicture}
    \useasboundingbox (-5,-5) rectangle (5,5);
    \ltdtsetobject[
        name={return "objs"},
        object={
            local emitter=Vector:new{3*cos(\i),3*sin(\i),0,1}
            local sink=Vector:new{0,0,0,1}
            local function field(u,v,w)
                local p=Vector:new{u,v,w,1}
                local ep=p:hsub(emitter)
                local sp=p:hsub(sink)
                local de=ep:hnorm()
                local ds=sp:hnorm()
                return ep:hscale(1/(de^3)):hsub(sp:hscale(1/(ds^3))):hnormalize()
            end
            local nu=\NU
            local nv=\NV
            local seed_radius=0.35
            local emitter_guard=0.31
            local sink_guard=0.31
            local step=0.2
            local max_steps=200
            local curves={}
            local n=1
            for i=0,nu-1 do
                local theta=tau*i/nu
                for j=0,nv-1 do
                    local phi=pi*(j+1)/(nv+1)
                    local direction=Vector.sphere(Vector:new{theta,phi,1})
                    local start=emitter:hadd(direction:hscale(seed_radius))
                    local alpha,intervals=Streamline.rk4(
                        field,
                        start,
                        step,
                        max_steps,
                        function(p)
                            return p:hdistance(emitter)<=emitter_guard
                                or p:hdistance(sink)<=sink_guard
                        end
                    )
                    curves[n]={
                        alpha=alpha,
                        fd=2/intervals
                    }
                    n=n+1
                end
            end
            return {
                emitter=emitter,
                sink=sink,
                curves=curves,
                view=Matrix.zyzrotation(Vector:new{7*pi/6,pi/3,pi/2})
            }
        }
    ]
    \ltdtappendlight[
        v={return Vector:new{1,1,1,1}}
    ]
    \ltdtappendsurface[
        uparams={return Vector:new{0,tau,16}},
        vparams={return Vector:new{0,pi,10}},
        v={
            return Vector.sphere(Vector:new{u,v,0.28}):hadd(objs.emitter)
        },
        transformation={return objs.view},
        fill options={
            return function(lighting)
                return "fill=red!"..tostring(floor(35+65*lighting)).."!black,draw=none"
            end
        }
    ]
    \ltdtappendsurface[
        uparams={return Vector:new{0,tau,16}},
        vparams={return Vector:new{0,pi,10}},
        v={
            return Vector.sphere(Vector:new{u,v,0.28}):hadd(objs.sink)
        },
        transformation={return objs.view},
        fill options={
            return function(lighting)
                return "fill=blue!"..tostring(floor(35+65*lighting)).."!black,draw=none"
            end
        }
    ]
    \foreach \j in {1,...,\streamcount} {
        \ltdtappendsurface[
            uparams={
                local c=objs.curves[\j]
                return Vector:new{c.fd,1-c.fd,40}
            },
            vparams={return Vector:new{0,tau,6}},
            v={
                local c=objs.curves[\j]
                local frame=Geometry.frenet(c.alpha,u,c.fd)
                return Vector:new{0,0.065*cos(v),0.065*sin(v),1}:multiply(frame)
            },
            transformation={return objs.view},
            fill options={
                return function(lighting)
                    return "fill=white!"..tostring(floor(20+80*lighting)).."!black,draw=none"
                end
            }
        ]
    }
    \ltdtdisplaysimplices
\end{tikzpicture}}
\end{document}